\section{Background and Related Work}
\label{sec:related-work}

This section reviews the evolution of quality assurance practices from traditional deterministic software systems to modern non-deterministic LLM-powered applications, and explores how self-testing practices have been adapted to serve as automated quality gates.

\subsection{Testing Traditional vs. LLM Applications}
Traditional software engineering relies heavily on deterministic testing methodologies, such as unit, integration, and end-to-end (E2E) testing \cite{wohlin2012experimentation}. These approaches assume that given a specific input, the system will consistently produce a specific, predictable output. However, testing LLM applications introduces a paradigm shift due to their probabilistic nature. Dobslaw and Feldt \cite{dobslaw2025testing} present a faceted taxonomy of challenges in testing LLM-based software, emphasizing that LLMs are prone to hallucinations, reasoning errors, and context sensitivity, making strict deterministic assertions nearly impossible.

Recent efforts in behavioral testing for NLP models, such as the CheckList framework by Ribeiro et al.~\cite{ribeiro2020beyond}, introduced the concept of testing linguistic capabilities via input perturbations. While CheckList works well for isolated NLP tasks (e.g., sentiment analysis), production LLM applications---like multi-agent orchestrators \cite{wu2024autogen, hong2024metagpt}---require evaluating the entire conversational pipeline, including context retrieval, tool use, and safety alignment.

Holistic benchmarking suites such as HELM \cite{liang2023helm}, MMLU \cite{hendrycks2021mmlu}, and BIG-bench \cite{srivastava2023bigbench} have advanced the field by evaluating models across diverse tasks. MT-Bench and Chatbot Arena \cite{zheng2023judging} specifically address LLM-judge evaluation for open-ended conversation. However, these benchmarks focus on model-level capabilities rather than application-level quality gates for production deployment.

\subsection{Automated Evaluation of LLMs}
To overcome the limitations of human evaluation, which is costly and difficult to scale~\cite{nushi2017human}, the industry has adopted automated evaluation techniques. A prominent approach is ``LLM-as-a-Judge,'' where a capable reasoning model evaluates the outputs of the system under test against a rubric \cite{zheng2023judging}. While this accelerates iteration cycles, Shankar et al. \cite{shankar2024who} argue that automated evaluators themselves require rigorous alignment and can introduce systemic biases.

The ML Test Score rubric proposed by Breck et al.~\cite{breck2017mltest} provides a framework for assessing ML production readiness, but focuses on traditional ML systems rather than the unique challenges of generative AI. Sculley et al. \cite{sculley2015debt} highlight the hidden technical debt in ML systems, including monitoring and testing debt, which becomes even more pronounced in LLM applications. Amershi et al. \cite{amershi2019se4ml} document nine stages of the ML workflow at Microsoft, noting that testing and monitoring remain the least mature practices.

In contrast to isolated benchmarking, evaluating production systems requires addressing the real-world usage distribution. Our framework builds upon automated evaluation but integrates it directly into the release pipeline as a strict quality gate, rather than an offline benchmarking tool.
This emphasis is complementary to broader release-readiness harnesses that combine evaluation, observability, and CI gating across heterogeneous LLM/RAG scenarios \cite{maiorano2026readiness}. Relative to that broader framing, the present work narrows the scope to automated self-testing as a longitudinal release-control mechanism for a single deployed multi-agent system.

\subsection{Continuous Delivery and Quality Gates}
Quality gates serve as decision checkpoints in software delivery pipelines~\cite{schermann2016qualitygates, shahin2017cicd}. Shahin et al. \cite{shahin2017cicd} provide a comprehensive review of continuous integration, delivery, and deployment practices, noting that automated quality gates are essential for maintaining release velocity without sacrificing quality. Metamorphic testing \cite{segura2016metamorphic} offers a complementary approach for systems where oracles are difficult to define---a property shared by LLM applications.

\subsection{Safety and Alignment of LLM Systems}
The safety dimension of our quality gate framework draws on recent work in LLM safety. Red teaming methodologies~\cite{ganguli2022redteaming, perez2022redteaming} have established systematic approaches for discovering harmful behaviors. Constitutional AI \cite{bai2022constitutional} and RLHF \cite{ouyang2022instructgpt} provide training-time alignment, but production systems still require runtime safety evaluation. Greshake et al. \cite{greshake2023injection} demonstrate that indirect prompt injection poses a significant threat to LLM-integrated applications, motivating the adversarial test tier in our question bank.

\subsection{Dogfooding in Software Engineering}
``Dogfooding''---the practice of an organization using its own product---has long been an informal mechanism for discovering usability issues and bugs before public release~\cite{harrison2006dogfooding}. Feitelson et al. \cite{feitelson2013facebook} describe how Facebook's deployment process relies on internal staging environments (H1/H2 servers) for gradual rollout, while Rossi et al. \cite{rossi2016facebook} explicitly document that ``dog-fooding and obtaining feedback from alpha and beta customers is critical to maintaining release quality'' in their study of Facebook's mobile deployment pipeline.

In traditional contexts, dogfooding is largely treated as a qualitative beta-testing phase. For LLM applications, automated self-testing provides a continuous stream of realistic conversational traces that can be codified into regression tests. Our work formalizes this practice into a quantifiable \textit{Quality Gate} framework following established case study methodology \cite{runeson2009guidelines, basili1994gqm}, closing the gap between internal testing and automated release decisions.

\subsection{Concurrent Automated Testing Frameworks}
\label{sec:concurrent}

Two concurrent works address related but distinct aspects of LLM application testing. STELLAR \cite{sorokin2026stellar} proposes a search-based framework that systematically generates failure-inducing inputs by discretizing linguistic and semantic feature dimensions (content, style, and perturbation). Applied to safety-critical (SafeQA) and retrieval-augmented (NaviQA) settings, STELLAR consistently uncovers more failing inputs than random, combinatorial, or coverage-based baselines. However, STELLAR's primary contribution is \emph{test input generation and failure discovery}; it does not formalize multi-dimensional quality gates or deterministic release decisions, and its evaluation is conducted offline rather than embedded in a CI/CD pipeline.

Multi-Agent LLM Committees \cite{karanam2025committees} orchestrate a committee of LLM agents using a three-round voting protocol to achieve consensus in autonomous beta testing. Evaluated across 84 runs with 9 distinct testing personas, the committee achieves an F1 score of 0.91 for regression detection versus 0.78 for single-agent baselines. Their focus is \emph{collaborative bug detection via committee consensus}, not the operationalization of release readiness across orthogonal quality dimensions such as evidence coverage or safety pass rate.

Our framework is distinguished from both by (i) formalizing five empirically grounded quality dimensions into a deterministic PROMOTE/HOLD/ROLLBACK gate, (ii) embedding the gate natively in a CI/CD workflow triggered on every merge, and (iii) empirically demonstrating the unique role of evidence coverage as a severe-regression discriminator across 38 longitudinal runs---a construct neither concurrent work addresses. Table~\ref{tab:concurrent_comparison} summarizes the key differences.

\begin{table}[htbp]
\renewcommand{\arraystretch}{1.2}
\centering
\footnotesize
\setlength{\tabcolsep}{3pt}
\caption{Qualitative comparison with concurrent automated testing frameworks for LLM applications (STELLAR~\cite{sorokin2026stellar}; MA-Committees~\cite{karanam2025committees}).}
\label{tab:concurrent_comparison}
\resizebox{\columnwidth}{!}{%
\begin{tabular}{p{2.0cm}p{1.7cm}p{1.7cm}p{1.9cm}}
\toprule
\textbf{Aspect} & \textbf{Ours} & \textbf{STELLAR} & \textbf{MA-Comm.} \\
\midrule
Primary goal & Release gating & Failure discovery & Bug detection \\
Quality dims. & 5D gate & Linguistic feats. & Persona diversity \\
Release decision & P/H/R & Failure set & Consensus vote \\
Evidence cov. & Formal dim. & Not addressed & Not addressed \\
Safety as dim. & \checkmark & Partial (SafeQA) & Adversarial persona \\
CI/CD & Native & Offline & Offline \\
Study scale & 38 runs & 2 use cases & 84 runs \\
Thresholds & Empirical & Search fitness & Vote threshold \\
\bottomrule
\end{tabular}}
\setlength{\tabcolsep}{6pt}
\end{table}


In the next section, we detail the architecture of this framework and the specific quality dimensions employed.

